\begin{lemma}[Spatial regularity of the forward contact envelope]
Let $F$ be a finite-dimensional real inner-product space, let
$1\le C_0<\infty$ and $0\le\kappa<\infty$, and let $v\in\Aclass_{m,C_0}$ in the
sense of \noderef{IsFitClass} be centered at $(0,1)$ with constant $C_0$ in the
sense of \noderef{IsCenteredAt} and satisfy the horizontal Jensen bound with
constant $\kappa$ at every time in the sense of \noderef{HasJensenBound}. Let
$f=W=\Cenv_tv-\kappa^2t/2$ be the contact envelope of \noderef{ContactEnv}, and
write $f^\ell=\Penv_\ell f$ for the forward envelope of \noderef{QuadEnvP}. Then
for every $t,\ell>0$ and every $x\in F$ the optimization defining
$f^\ell(x,t)$ has a unique maximizer $y=y(x,t,\ell)$, the function
$f^\ell(\cdot,t)$ is differentiable at $x$ with
\begin{equation}\label{eq:dp-gradient}
  \nabla_x f^\ell(x,t)=\frac{y-x}{\ell},
\end{equation}
and for every $h\in F$ the two supporting inequalities hold:
\begin{equation}\label{eq:dp-supports}
\begin{split}
  f^\ell(x,t)+\nabla_x f^\ell(x,t)\cdot h-\frac{|h|^2}{2\ell}
  &\le f^\ell(x+h,t)\\
  &\le f^\ell(x,t)+\nabla_x f^\ell(x,t)\cdot h+\frac{|h|^2}{2t}.
\end{split}
\end{equation}
In particular the spatial Hessian of $f^\ell(\cdot,t)$ is bounded below by
$-\ell^{-1}I$ in the semiconvexity sense and above by $t^{-1}I$ in the
semiconcavity sense.
\end{lemma}
\begin{proof}
Write $f=W=\Cenv_tv-\kappa^2t/2$ for the contact envelope of
\noderef{ContactEnv}, fix $t,l>0$ and $x\in F$, and put $s=t+l>0$.

\emph{Step 0: regularity of the integrand on positive times.}
The hypotheses of \noderef{ContactEnvPositiveTimeRegularity} are exactly those assumed
here: $1\le C_0$, $0\le\kappa$, $v$ in the class $\Aclass_{m,C_0}$ of
\noderef{IsFitClass}, centered at $(0,1)$ with constant $C_0$ in the sense of
\noderef{IsCenteredAt}, and satisfying the horizontal Jensen bound with constant $\kappa$
at every time in the sense of \noderef{HasJensenBound}. That lemma therefore gives three
facts about $f$, which are the only properties of $f$ used below:
\begin{enumerate}
\item $f$ is continuous on $\{z:z_t>0\}$;
\item $f$ has subquadratic spatial growth on compact subintervals of $(0,\infty)$ in the
sense of \noderef{SubquadraticGrowth};
\item for every $\sigma>0$ the function $y\mapsto f(y,\sigma)-|y|^2/(2\sigma)$ is concave
on $F$.
\end{enumerate}

\emph{Step 1: $f(\cdot,s)-|\cdot|^2/(2s)$ is concave.}
Apply part (3) of Step 0 with $\sigma=s=t+l>0$:
\begin{equation}\label{eq:dpsr-concave}
  y\longmapsto f(y,s)-\frac{|y|^2}{2s}
\end{equation}
is concave on $F$.

\emph{Step 2: strict concavity of the objective, hence uniqueness.}
Write, for $y\in F$,
\[
  \Xi(y)=f(y,s)-\frac{|y-x|^2}{2l}
   =\Bigl(f(y,s)-\frac{|y|^2}{2s}\Bigr)
     +\Bigl(\frac{|y|^2}{2s}-\frac{|y-x|^2}{2l}\Bigr).
\]
The first bracket is concave by \eqref{eq:dpsr-concave}. The second is a
quadratic function of $y$ whose quadratic part is
$\tfrac12\bigl(\tfrac1s-\tfrac1l\bigr)|y|^2=-\tfrac{t}{2ls}|y|^2$, strictly
concave because $t,l,s>0$. Hence $\Xi$ is strictly concave: for $y_0\ne y_1$ and
$\theta\in(0,1)$,
$\Xi((1-\theta)y_0+\theta y_1)>(1-\theta)\Xi(y_0)+\theta\Xi(y_1)$. In
particular $\Xi$ has at most one maximizer.

\emph{Step 3: existence of the maximizer.}
Apply \noderef{QuadEnvAttained} to $f$ with $I=(0,\infty)$: its hypotheses are parts (1)
and (2) of Step 0. Part (1), at the center $(x,t)$ and the length $l>0$, whose time
arguments $t$ and $t+l=s$ both lie in $(0,\infty)$, gives $y\in F$ with
\begin{equation}\label{eq:dpsr-attain}
  \Penv_lf(x,t)=f(y,s)-\frac{|y-x|^2}{2l}=\Xi(y)=\max_{y'\in F}\Xi(y') .
\end{equation}
By Step 2 this maximizer is unique, which is the first two assertions of the
lemma.

The hypotheses of part (1) just applied constrain only the two time arguments,
namely $t$ and $t+l=s$, and both lie in $I=(0,\infty)$ independently of the
spatial center. So the identical application of part (1) at the center
$(x',t)$, for an arbitrary $x'\in F$ and the same length $l>0$, gives a point
$y'\in F$ at which the supremum defining $\Penv_lf(x',t)$ in
\noderef{QuadEnvP} is attained:
\begin{equation}\label{eq:dpsr-attain-any}
  \Penv_lf(x',t)=f(y',s)-\frac{|y'-x'|^2}{2l}
   =\max_{w\in F}\Bigl\{f(w,s)-\frac{|w-x'|^2}{2l}\Bigr\}
  \qquad(x'\in F).
\end{equation}
Two consequences are recorded here once and used below without further comment.
First, $\Penv_lf(x',t)$ is a real number equal to a maximum at every center, so
it is a genuine least upper bound of the competitor values there and not the
junk value fixed by the preamble convention for a supremum over an empty or
unbounded-above set. Second, it therefore dominates the value of each single
competitor:
\begin{equation}\label{eq:dpsr-competitor}
  \Penv_lf(x',t)\ge f(w,s)-\frac{|w-x'|^2}{2l}\qquad(x',w\in F).
\end{equation}

\emph{Step 4: the lower supporting inequality.}
Let $h\in F$. Keeping the maximizer $y$ of \eqref{eq:dpsr-attain} fixed and
using it as a competitor at the center $x+h$, \eqref{eq:dpsr-competitor} with
$x'=x+h$ and $w=y$ gives
\[
  \Penv_lf(x+h,t)\ge f(y,s)-\frac{|y-x-h|^2}{2l}
   =\Penv_lf(x,t)+\frac{|y-x|^2-|y-x-h|^2}{2l},
\]
the equality by \eqref{eq:dpsr-attain}. Since
$|y-x-h|^2=|y-x|^2-2\langle y-x,h\rangle+|h|^2$, the fraction equals
$\langle(y-x)/l,h\rangle-|h|^2/(2l)$, so
\begin{equation}\label{eq:dp-lower-support}
  \Penv_lf(x+h,t)\ge\Penv_lf(x,t)
    +\Bigl\langle\frac{y-x}{l},h\Bigr\rangle-\frac{|h|^2}{2l}.
\end{equation}

\emph{Step 5: the upper supporting inequality.}
The three hypotheses of \noderef{QuadEnvPConcaveCorrection}, for this $f$ and these
$t,l>0$, are available: its hypothesis (1) is \eqref{eq:dpsr-concave}; its hypothesis (2),
attainment of the defining supremum at every center, is
\eqref{eq:dpsr-attain-any}; and its hypothesis (3), domination of every competitor at
every center, is \eqref{eq:dpsr-competitor}. That lemma therefore gives that
\[
  \Lambda(x'):=\Penv_lf(x',t)-\frac{|x'|^2}{2t}
\]
is concave on $F$. It is real-valued at every $x'\in F$ by
\eqref{eq:dpsr-attain-any}, which exhibits $\Penv_lf(x',t)$ as a maximum of real
numbers.

Since $\Lambda$ is real-valued and concave on $F$, the function $\psi:=-\Lambda$
is real-valued and convex on $F$; and $F$ is an open subset of $F$ containing
$x$. So \noderef{ConvexSubgradientExists}, applied with the open set
$\Omega=F$, the function $\psi=-\Lambda$ and the point $z=x\in F$, provides
$q_0\in F$ with
\[
  -\Lambda(x)+\langle q_0,\chi-x\rangle\le-\Lambda(\chi)\qquad(\chi\in F).
\]
Multiplying by $-1$ and writing $q:=-q_0$ gives
\[
  \Lambda(\chi)\le\Lambda(x)+\langle q,\chi-x\rangle\qquad(\chi\in F).
\]
Taking $\chi=x+h$ and adding
$|x+h|^2/(2t)=|x|^2/(2t)+\langle x,h\rangle/t+|h|^2/(2t)$ to both sides,
\begin{equation}\label{eq:dpsr-upper}
  \Penv_lf(x+h,t)\le\Penv_lf(x,t)+\langle q',h\rangle+\frac{|h|^2}{2t},
  \qquad q':=q+\frac xt .
\end{equation}

\emph{Step 6: identification of $q'$ and differentiability.}
Let $e\in F$ be a unit vector and $\rho>0$. Applying
\eqref{eq:dp-lower-support} and \eqref{eq:dpsr-upper} with $h=\rho e$,
\[
  \rho\Bigl\langle\frac{y-x}{l},e\Bigr\rangle-\frac{\rho^2}{2l}
   \le\Penv_lf(x+\rho e,t)-\Penv_lf(x,t)
   \le\rho\langle q',e\rangle+\frac{\rho^2}{2t},
\]
so $\bigl\langle(y-x)/l-q',e\bigr\rangle\le\rho\bigl(\tfrac1{2t}+\tfrac1{2l}\bigr)$;
letting $\rho\downarrow0$ gives
$\bigl\langle(y-x)/l-q',e\bigr\rangle\le0$. As $e$ ranges over all unit
vectors, this forces $q'=(y-x)/l$. Hence
\eqref{eq:dp-lower-support} and \eqref{eq:dpsr-upper} read
\[
  -\frac{|h|^2}{2l}
   \le\Penv_lf(x+h,t)-\Penv_lf(x,t)-\Bigl\langle\frac{y-x}{l},h\Bigr\rangle
   \le\frac{|h|^2}{2t}\qquad(h\in F),
\]
which are the two supporting inequalities \eqref{eq:dp-supports}, and which show
that $x'\mapsto\Penv_lf(x',t)$ is differentiable at $x$ with gradient
$(y-x)/l$, since the middle term is $O(|h|^2)=o(|h|)$. This is
\eqref{eq:dp-gradient}.
\end{proof}
